% Recuperación documental para el artículo V; RESULTADO_RECUPERADO.
% Propietario: /Users/ruben/Documents/New project/output/REVISION_EDITORIAL_HMT_MD_20260905/01_FUENTES/INTEGRAL/tree/New project/output/ENSAMBLADO_CAUSAL_RESTAURADO_HMT_MD_20260905/fuente/manuscrito/sections/md/delta_127/04b_completaciones_cuanticas.tex
% Líneas de origen: 1--334.
% REV02: construcción algebraica; la aplicación estadística se reúne en 40.
% Correspondencia editorial: gestion/CONSOLIDACION_PAULI_FOCK_REV02.json.
\section{Construcción algebraica de las completaciones cuánticas sobre APP}
\label{v:rec:chap:completaciones-cuanticas}
\label{v:sec:completaciones-algebraicas}

APP y TPK proporcionan estados orientados, transporte y paridad de hoja.
Su realización lineal permite construir las álgebras simétrica y exterior,
con sus operadores de creación, aniquilación y ocupación. Esta sección
establece esas construcciones algebraicas, el sector celular de campos y
el acoplamiento reversible de medición. El transporte de las completaciones
por refinamiento se demuestra en la
sección~\ref{v:rec:fock_gibbs:section:01}; la aplicación de sus espectros de
ocupación a los recuentos y al equilibrio se desarrolla en las
secciones~\ref{v:rec:chap:ocupacion-estadistica}--\ref{v:rec:ocupacion:section:03}.
Cada tramo conserva las condiciones de la representación que utiliza.

\subsection{Espacio lineal de estados asociado a una realización de partícula}
\label{v:rec:completaciones:section:02}

Sea \(S\) un conjunto finito de clases de rutas TPK a la profundidad elegida
y sea
\[
 \mathcal H_1=\ell^2(S)
\]
el espacio de Hilbert con base ortonormal \((e_s)_{s\in S}\). Una elección de
pesos positivos \(E_s\) define el operador diagonal
\[
 He_s=E_se_s.
\]
La elección de \(S\), de su cociente por simetrías y de los pesos forma parte
del modelo. Una vez fijada, las completaciones multipartícula son canónicas.

\subsection{Completaciones simétrica y exterior}
\label{v:rec:completaciones:section:03}

Definimos
\[
 \mathcal F_+(\mathcal H_1)=
 \bigoplus_{n\geq0}\operatorname{Sym}^n\mathcal H_1,
 \qquad
 \mathcal F_-(\mathcal H_1)=
 \bigoplus_{n\geq0}\bigwedge^n\mathcal H_1.
\]
La primera es la completación bosónica; la segunda, la fermiónica. En ambas,
el sumando de grado cero es el espacio del vacío. Las potencias exteriores
se dotan del producto escalar para el cual las cuñas ordenadas de vectores
de una base ortonormal son ortonormales; las potencias simétricas se utilizan
con la normalización usual de los operadores de creación y aniquilación.

\subsection{Graduación de hoja y álgebra de intercambio}
\label{v:rec:completaciones:section:04}

Supongamos que la hoja orientada define un carácter aditivo
\[
 p:\operatorname{Mor}(\mathsf P_{\rm TPK})\longrightarrow\mathbb Z/2\mathbb Z,
 \qquad p(\gamma\eta)=p(\gamma)+p(\eta).
\]
Sobre el espacio lineal graduado correspondiente se introduce la simetría
\[
 \tau(v\otimes w)=(-1)^{p(v)p(w)}w\otimes v.
\]

\begin{theorem}[Completación estadística determinada por la paridad]
\label{v:rec:completaciones:statement:02}
El cociente del álgebra tensorial por las relaciones
\[
 v\otimes w-(-1)^{p(v)p(w)}w\otimes v
\]
es simétrico sobre el sector par y exterior sobre el sector impar. Dos estados
impares idénticos tienen producto nulo; los estados pares admiten ocupación
arbitraria.
\end{theorem}

\begin{proof}
Si \(p(v)=p(w)=0\), la relación es \(vw=wv\). Si
\(p(v)=p(w)=1\), es \(vw=-wv\); tomando \(v=w\) y trabajando en
característica cero se obtiene \(v^2=0\). Los sectores mixtos conmutan sin
signo adicional. Esta es la propiedad universal del álgebra simétrica
graduada.
\end{proof}

\subsection{Operadores canónicos e idempotencia de la ocupación}
\label{v:pauli:car-construccion}

El carácter de paridad determina la completación; el producto de esa
completación determina a continuación el álgebra de operadores. En el
sector exterior, la creación \(a_s^\dagger\) es el producto exterior por
\(e_s\), y \(a_s\) es su adjunto.

\begin{theorem}[Principio de exclusión de Pauli en la completación exterior]
\label{v:rec:completaciones:statement:01}
\label{v:rec:thm:principio-exclusion-pauli}
Para todo modo \(s\in S\), los operadores de creación y aniquilación de la
completación exterior satisfacen
\begin{equation}
 \{a_s,a_t^\dagger\}=\delta_{st}I,
 \qquad
 \{a_s,a_t\}=\{a_s^\dagger,a_t^\dagger\}=0.
 \label{v:pauli:eq:car}
\end{equation}
En particular, el operador de ocupación
\(N_s=a_s^\dagger a_s\) es un proyector:
\begin{equation}
 N_s^2=N_s,
 \qquad \Spec(N_s)\subseteq\{0,1\}.
 \label{v:pauli:eq:proyeccion-ocupacion}
\end{equation}
Los autovalores \(0\) y \(1\) corresponden, respectivamente, a la ausencia
y a la presencia del modo \(s\).
\end{theorem}

\begin{proof}
Sobre una cuña ordenada, \(a_s\) suprime el factor \(e_s\), cuando está
presente, con el signo de su posición; si está ausente, da cero. Componer
esta operación con el producto exterior demuestra las relaciones
\eqref{v:pauli:eq:car}: para \(s=t\), las dos composiciones cubren los casos
complementarios de ausencia y presencia; para \(s\ne t\), los signos de
intercambio se cancelan. La antisimetría implica
\(e_s\wedge e_s=0\) y
\(a_s^2=(a_s^\dagger)^2=0\). Entonces
\begin{equation}
 N_s^2=a_s^\dagger a_s a_s^\dagger a_s
 =a_s^\dagger(I-a_s^\dagger a_s)a_s
 =a_s^\dagger a_s=N_s.
 \label{v:pauli:eq:prueba-idempotencia}
\end{equation}
El operador es autoadjunto, y todo autovalor de un idempotente satisface
\(\lambda^2=\lambda\). Por tanto, su espectro está contenido en
\(\{0,1\}\). El vacío y el estado \(e_s\) realizan ambos valores
\cite{Pauli1925,Dirac1926}.
\end{proof}

En el sector simétrico, los operadores bosónicos \(b_s^\dagger,b_s\)
actúan sobre la base normalizada de ocupaciones por los factores
\(\sqrt{m_s+1}\) y \(\sqrt{m_s}\), respectivamente. De sus composiciones
se obtienen
\begin{equation}
 [b_s,b_t^\dagger]=\delta_{st}I,
 \qquad [b_s,b_t]=[b_s^\dagger,b_t^\dagger]=0.
 \label{v:pauli:eq:ccr}
\end{equation}
Estas identidades se entienden sobre el dominio algebraico de partículas
finitas, estable bajo los operadores considerados.

El teorema~\ref{v:rec:completaciones:statement:02} convierte una paridad composicional TPK en una regla exacta de
ocupación. Para obtener el teorema físico de espín--estadística hay que
construir además una teoría local lorentziana positiva y demostrar que la
paridad de hoja coincide con el espín bajo rotaciones. La igualdad no se
deduce únicamente de que ambos objetos tomen valores módulo dos.

\subsection{Holonomía de hoja y periodicidad de \texorpdfstring{\(2\pi/4\pi\)}{2pi/4pi}}
\label{v:rec:completaciones:section:05}

La paridad anterior admite una realización geométrica natural cuando una
extensión superviviente posee un retorno cerrado cuya cubierta de orientación
es no trivial. Sea \(\rho\) el generador del retorno y sea
\[
 \varepsilon:\langle\rho\rangle\longrightarrow\{-1,+1\}
\]
su carácter de holonomía, con \(\varepsilon(\rho)=-1\). El fibrado real
asociado es entonces la línea de orientación de una banda de Möbius: una
vuelta cierra la proyección geométrica, pero cambia la hoja levantada; dos
vueltas restauran la orientación.

Para escribir esta propiedad angularmente, fijemos una representación
unitaria \(U\) de la cubierta de rotaciones sobre el espacio graduado
\(\mathcal H_1=\mathcal H_{\bar0}\oplus\mathcal H_{\bar1}\), compatible con
la paridad de hoja en el sentido
\[
 U(2\pi)v=(-1)^{p(v)}v
\]
para todo vector homogéneo \(v\). La compatibilidad es una condición de la
realización: identifica el generador de la cubierta de orientación con el
carácter \(p\) ya construido sobre las rutas.

\begin{theorem}[Restauración de orientación y álgebra de intercambio]
\label{v:rec:thm:dos-cuatro-pi-estadistica}
Bajo la compatibilidad precedente,
\[
 U(2\pi)|_{\mathcal H_{\bar0}}=I,
 \qquad
 U(2\pi)|_{\mathcal H_{\bar1}}=-I,
 \qquad
 U(4\pi)=I.
\]
Además, el mismo carácter \(p\) determina la simetría graduada
\[
 \tau(v\otimes w)=(-1)^{p(v)p(w)}w\otimes v.
\]
Por tanto, el sector de holonomía par se completa mediante potencias
simétricas y el sector de holonomía impar mediante potencias exteriores. En
el segundo se cumple la exclusión \(v\wedge v=0\).
\end{theorem}

\begin{proof}
La primera y la segunda igualdad son la condición de compatibilidad evaluada
en cada componente homogénea. Al iterar una vuelta,
\[
 U(4\pi)v=U(2\pi)^2v=(-1)^{2p(v)}v=v,
\]
lo que demuestra la restauración. La afirmación sobre el intercambio es el
teorema~\ref{v:rec:completaciones:statement:02}: para dos vectores
pares el signo es positivo, mientras que para dos impares es negativo.
\end{proof}

Conviene distinguir tres retornos. El cierre angular tras \(2\pi\) es un
retorno de la proyección; el cierre tras \(4\pi\) restaura la orientación
levantada; ninguno obliga por sí solo a que todas las coordenadas del registro
regresen a su valor inicial. Esta separación es la análoga dimensional de la
distinción HMT entre retorno de fase y retorno de estado.

\subsubsection{Las hojas suma y producto como soportes de las \(2\) completaciones}
\label{v:rec:completaciones:section:06}

La tabla APP contiene dos evaluaciones sobre un mismo soporte. La hoja
aditiva actúa por traslaciones y conserva la multiplicidad de las
continuaciones; en este sentido, acumula. La hoja multiplicativa es
permutativa sobre las unidades y contrae fibras sobre los no invertibles; en
este sentido, pliega. Esta incompatibilidad es anterior a la elección
estadística y explica por qué TPK debe conservar dos genealogías en lugar de
reducirlas a una sola palabra visible.

La asignación estadística no se obtiene, sin embargo, sustituyendo
verbalmente «suma» por «bosón» y «producto» por «fermión». El mapa exacto es
el carácter de paridad
\[
 p:\operatorname{Mor}(\mathsf P_{\rm TPK})\longrightarrow\mathbb Z/2\mathbb Z.
\]
Cuando el transporte de una extensión cae en el sector \(p=0\), su
completación admite ocupación acumulativa y es simétrica. Cuando cae en
\(p=1\), la reversión de hoja introduce el signo de intercambio y la
completación es exterior. La suma y el producto suministran las dos
geometrías locales; el registro, la orientación y el carácter \(p\) deciden en
qué sector se realiza una ruta.

La periodicidad \(4\pi\) y la exclusión fermiónica quedan coordinadas por
el mismo carácter de hoja. Sus consecuencias para los recuentos finitos y
las ocupaciones de equilibrio se obtienen, con el Hamiltoniano y las
condiciones térmicas explícitos, en la
sección~\ref{v:rec:chap:ocupacion-estadistica} y en el
teorema~\ref{v:rec:thm:factorizacion-gran-canonica-rev3}.

\subsection{El complejo celular del toro APP}
\label{v:rec:completaciones:section:08}

El grafo de APP es el uno-esqueleto de una descomposición celular del toro
\(T^2\) con \(9\times9\) vértices. Sea
\[
 C^0\xrightarrow{d_0}C^1\xrightarrow{d_1}C^2
\]
su complejo de cocadenas reales. Elegido el producto escalar uniforme,
definimos la codiferencial \(\delta=d^*\) y el laplaciano
\[
 \Delta_k=d_{k-1}\delta_k+\delta_{k+1}d_k.
\]

\begin{theorem}[Sector armónico del toro digital]
\label{v:rec:completaciones:statement:04}
La dimensión del espacio de uno-formas armónicas es dos:
\[
 \dim\ker\Delta_1=b_1(T^2)=2.
\]
Puede elegirse una base representada por los flujos uniformes horizontal y
vertical. Ambos son cerrados, cocerrados y no son gradientes de funciones
globales periódicas.
\end{theorem}

\begin{proof}
El teorema de Hodge para complejos celulares finitos identifica
\(\ker\Delta_1\) con \(H^1(T^2;\mathbb R)\). La cohomología del toro tiene
rango dos. Directamente, los flujos constantes horizontal y vertical tienen
divergencia y rotacional nulos; sus integrales sobre los dos ciclos
fundamentales son independientes, por lo que no son exactos.
\end{proof}

El espacio de conexiones módulo transformaciones de calibre es
\[
 C^1/\operatorname{im}d_0.
\]
En un complejo finito provisto del producto escalar anterior, cada órbita
admite un representante de Coulomb en \(\ker\delta_1\); dentro de la órbita,
la elección es única módulo el estabilizador \(\ker d_0\) del parámetro de
calibre. El cociente cohomológico correcto es
\[
 \ker d_1/\operatorname{im}d_0
 \simeq \ker\Delta_1 .
\]
Sus dos clases armónicas constituyen un candidato matemático natural para
dos modos globales libres. No son todavía las dos polarizaciones locales de
un fotón en espacio-tiempo \(3+1\): para esa identificación se necesita un
límite dinámico, una relación de dispersión y un grupo físico de calibre.

\subsection{Ley de Gauss discreta}
\label{v:rec:completaciones:section:09}

Para una 1-cocadena orientada \(F\) y un conjunto de vértices \(V\), la
divergencia se define por
\[
 (\operatorname{div}F)(v)=
 \sum_{e:v\to w}F(e)-\sum_{e:w\to v}F(e).
\]
Al sumar sobre \(V\), toda arista interior aparece una vez con cada signo y
se cancela. Por tanto,
\[
 \boxed{
 \sum_{v\in V}\operatorname{div}F(v)
 =\operatorname{Flux}_{\partial V}(F).}
\]
Esta identidad es una ley de Gauss exacta en el grafo. Interpretar
\(\operatorname{div}F\) como carga eléctrica exige el carácter de carga y la
realización dimensional construidos en Mecánica Dimensional.

\subsection{Aparato de medición reversible}
\label{v:rec:completaciones:section:10}

Sea \(S\) un conjunto de estados y sea \(r:S\to\mathbb Z/9\mathbb Z\) una
lectura. El acoplamiento a un registro nonádico se define por
\[
 U_r:S\times\mathbb Z/9\mathbb Z\longrightarrow
 S\times\mathbb Z/9\mathbb Z,
 \qquad
 U_r(s,a)=(s,a+r(s)).
\]

\begin{proposition}[Reversibilidad del acoplamiento]
\label{v:rec:completaciones:statement:05}
La aplicación \(U_r\) es una permutación y
\[
 U_r^{-1}(s,a)=(s,a-r(s)).
\]
La pérdida de información sólo aparece después de aplicar un lector al
registro o de condicionar el estado por el resultado.
\end{proposition}

\begin{proof}
La fórmula propuesta para la inversa satisface ambas composiciones por la
ley de grupo de \(\mathbb Z/9\mathbb Z\).
\end{proof}

Esta proposición formaliza la definición operativa de medición: acoplar,
leer y restringir futuros compatibles son tres aplicaciones diferentes. El
acoplamiento completo es reversible; el cociente observacional puede no serlo.

\subsection{Correspondencias algebraicas y realizaciones físicas}
\label{v:rec:completaciones:section:11}

Las completaciones de Fock, la exclusión condicionada a paridad, el sector
armónico de APP, Gauss discreto y la reversibilidad del aparato son teoremas
matemáticos completos. Sus identificaciones físicas forman un diagrama que
debe conmutar:
\[
 \text{paridad de hoja}
 \longrightarrow\text{intercambio graduado}
 \longrightarrow\text{CAR/CCR}
 \longrightarrow\text{estadística física},
\]
\[
 H^1(\mathcal A_9)\simeq\ker\Delta_1
 \hookrightarrow C^1/\operatorname{im}d_0
 \longrightarrow\text{realización fotónica}.
\]
En la primera línea, las dos primeras flechas poseen construcciones
explícitas; la última requiere la dinámica y las condiciones de equilibrio
correspondientes. En la segunda, la identificación de Hodge y la inclusión
de las clases armónicas en el espacio de conexiones módulo calibre son las
construcciones de la sección~\ref{v:rec:completaciones:section:08}. El autovalor
cero es el del laplaciano celular: identificar estos modos con excitaciones
sin masa y con las dos polarizaciones locales del fotón requiere además la
dinámica, la relación de dispersión y la realización física de calibre allí
indicadas. Esta separación permite desarrollar los puentes sin confundir una
analogía de cardinalidad con una representación.
